\documentclass{article}

\usepackage{amsmath,amssymb}
\usepackage{xurl}
\usepackage[hidelinks]{hyperref}
\setlength{\emergencystretch}{2em}

% Shared commands for notes in docs/paper-gaps/.

\DeclareUnicodeCharacter{2080}{\ifmmode{}_{0}\else\textsubscript{0}\fi}
\DeclareUnicodeCharacter{2081}{\ifmmode{}_{1}\else\textsubscript{1}\fi}
\DeclareUnicodeCharacter{2082}{\ifmmode{}_{2}\else\textsubscript{2}\fi}
\DeclareUnicodeCharacter{2099}{\ifmmode{}_{n}\else\textsubscript{n}\fi}

\newcommand{\Lean}{\textsc{Lean}}
\ExplSyntaxOn
\NewDocumentCommand{\leanid}{m}{
  \ifmmode
    \textsf{\detokenize{#1}}
  \else
    \group_begin:
    \urlstyle{sf}
    \tl_set:Nn \l_tmpa_tl {#1}
    \tl_replace_all:Nnn \l_tmpa_tl {\_} {_}
    \exp_args:NV \path \l_tmpa_tl
    \group_end:
  \fi
}
\ExplSyntaxOff
\providecommand{\tr}{\operatorname{tr}}
\newcommand{\tnleanIssueBase}{https://github.com/LionSR/TNLean/issues/}
\newcommand{\tnleanPrBase}{https://github.com/LionSR/TNLean/pull/}
\newcommand{\tnleanDocBase}{https://sirui-lu.com/TNLean/docs/}
\newcommand{\ghissue}[1]{\href{\tnleanIssueBase#1}{GitHub issue \##1}}
\newcommand{\ghpr}[1]{\href{\tnleanPrBase#1}{PR \##1}}
\newcommand{\doclink}[1]{\href{\tnleanDocBase#1.html}{\path{#1}}}

% Dirac notation and the identity operator, as in blueprint/src/macros/common.tex.
% \providecommand keeps note-local definitions authoritative.
\usepackage{dsfont}
\providecommand{\ket}[1]{|#1\rangle}
\providecommand{\bra}[1]{\langle #1|}
\providecommand{\braket}[2]{\langle #1 | #2 \rangle}
\providecommand{\ketbra}[2]{|#1\rangle\!\langle #2|}
\providecommand{\Id}{\mathds{1}}
\providecommand{\one}{\mathds{1}}
\providecommand{\C}{\mathbb{C}}
\providecommand{\MN}[1]{M_{#1}(\C)}
\providecommand{\MD}{M_D(\C)}

% Verdict marker: \gapnote{<kind>}{<status>}. Kind is the policy
% classification (clarification, local-correction, scope-restriction,
% unfaithful, false-source, open-gap); status is open, wip (elimination
% underway), resolved, or historical. Machine-read by the site index;
% prints a verdict line.
\newcommand{\gapnote}[2]{\par\noindent\textbf{Verdict:} #1 (#2).\par}


\title{Guide Properties in the Angular Separation of Dyadic Repaintings}
\author{}
\date{2026-10-07; resolved 2026-10-08}

\begin{document}
\maketitle
\gapnote{scope-restriction}{resolved}

\section{The source statement}

Section~7 of the September 24, 2026 polynomial PEPS manuscript\footnote{%
\url{https://github.com/openai/math/blob/adc7f1241b42e322a6451854ab7e4b4c146bf78a/preprints/Polynomial-PEPS-approximation-of-gapped-square-grid-ground-states-September-24-2026/build/main.tex},
file \texttt{06-geometry.tex}, lines 254--319.}
repaints an aligned square $S=(nr,ns)+[0,n]^2$ of the dyadic hierarchy from its
old label $A$ to its final label $B$. Along each edge it uses the normal ratio
$x=d/w(s)$ with $w(s)=10^{-3}\min(s,n-s)$, and it runs a central birth, seven
elementary births and deaths per edge, and one lens exchange per edge. Lemma~7.2
asserts that there is one constant $a_0>0$ such that, for every one of these
operations and for every point $y$ of the closure of its changed region,
\begin{align}
    \operatorname{dist}(y,\{f\ne P_\circ\})
        &\ge a_0\min\{n,d_{\mathcal V}(y)\},
    \label{eq:openai26_dyadic_geometry_guide_properties_clearance}
\end{align}
where $f$ is the surrounding guide of the schedule, $P_\circ$ the surrounding
label, and $d_{\mathcal V}$ the distance from the marked points; for the exchange
the obstructing set is the boundary of the lens $Y=\{-7<x<-7/2\}$.

\section{The earlier formal statement}

Initially the schedule of guides was not constructed. The
estimate~(\ref{eq:openai26_dyadic_geometry_guide_properties_clearance}) was
proved with $a_0=1/5000$ for every guide having the properties that the source's
proof uses:
\begin{itemize}
    \item for an edge birth or death, the guides before, after and surrounding
    the operation carry the source's normal words on the band $\{-8<x<2\}$, and
    the guides before and after agree off that band;
    \item for the central birth, the surrounding guide equals $A$ on the open
    square;
    \item for the lens exchange, the main guide equals $C$ on the exterior band
    $\{-8<x<0\}$, and the auxiliary guide carries the auxiliary word on
    $\{-8<x<2\}$ and equals $C$ off it.
\end{itemize}
The square is placed at the origin. These are the
declarations \leanid{TNLean.PEPS.Approximation.BandOperation.clearance},
\leanid{TNLean.PEPS.Approximation.centralBirth_clearance} and
\leanid{TNLean.PEPS.Approximation.lensExchange_clearance}.

The source derives the listed properties at lines 227--241: the bands
$\{-8<x<2\}$ along distinct edges are disjoint, so the operations along one edge
do not meet those along another, and the exterior part of a band stays in one
neighboring square, which carries a constant label. The disjointness is proved
as \leanid{TNLean.PEPS.Approximation.disjoint_edgeBand}. The formal statements
therefore cover the source's operations once the guides are built, and no
hypothesis absent from the source's argument is used. A translate of $S$ only
translates the coordinates, and sup distances are translation invariant.

\section{Resolution}

The schedule of guides is now constructed, and Lemma~7.2 is proved for it on
every square of the hierarchy, with $a_0=1/5000$.

The starting guide of a repainting of $S=[0,n]^2$ is any labelling of the plane
that is $A$ on the open square and carries one label $C_e$ on the open exterior
neighboring square across each edge~$e$. The guide after the central birth is
$B$ on the central region. The guide after the edges of a list $E$ have been
completed replaces the band $\{-8<x<2\}$ of each completed edge by the word
$C_e\mid B$ at $0$. During the construction along an edge~$e$, the main and
auxiliary guides are the completed guide and the uniform guide $C_e$ with the
band $\{-8<x<2\}$ of $e$ replaced by the current normal word. These are the
declarations \leanid{TNLean.PEPS.Approximation.RepaintingBaseline},
\leanid{TNLean.PEPS.Approximation.RepaintingBaseline.completedGuide},
\leanid{TNLean.PEPS.Approximation.RepaintingBaseline.mainGuide} and
\leanid{TNLean.PEPS.Approximation.RepaintingBaseline.auxGuide}.

The derivation of the guide properties at lines 227--241 is proved: by the
band disjointness, on the band of an edge not yet completed the completed guide
carries $C_e\mid A\mid B$ at $0,1$
(\leanid{TNLean.PEPS.Approximation.RepaintingBaseline.completedGuide_eq_edgeStartWord});
the lens exchange produces the main word
(\leanid{TNLean.PEPS.Approximation.RepaintingBaseline.lensExchange_mainGuide});
consecutive elementary operations on one sheet match; and after the four edges
the open square carries $B$ while every point outside the closed square keeps
its starting label. The estimates are transported to the translate
$(nr,ns)+[0,n]^2$ of the square, whose starting guide is uniform on $n$-blocks,
and stated in the source's form with closures:
\leanid{TNLean.PEPS.Approximation.blockRepainting_central},
\leanid{TNLean.PEPS.Approximation.blockRepainting_auxStep},
\leanid{TNLean.PEPS.Approximation.blockRepainting_mainStep} and
\leanid{TNLean.PEPS.Approximation.blockRepainting_exchange}. Their only
hypothesis is $n>0$; the block labels and the final label are arbitrary.

One convention remains, and it follows the source's own. A guide of the source
is specified by its open polygonal chambers, and its labels are read at the
lattice sites, the cell centers $(i+\tfrac12,j+\tfrac12)$, after one small
generic displacement (lines 62--80). A formal guide is a labelling of the whole
plane, and the estimates are proved for the labellings that the schedule
specifies at each operation; within one edge construction consecutive labellings
match exactly. Twice the schedule passes to a labelling that agrees with the one
reached only off finitely many curves. At the start of the construction along an
edge~$e$, the completed guide is replaced by the main guide with the starting
word $C_e\mid A\mid B$, which agrees with it off the curves $x=0$ and $x=1$ of
the band of $e$
(\leanid{TNLean.PEPS.Approximation.RepaintingBaseline.completedGuide_eq_mainGuide_start}).
At the end of a repainting, the final guide is replaced by the block guide of
the next repainting, which agrees with it off the boundary of the repainted
square (\leanid{TNLean.PEPS.Approximation.blockRepainting_final}).

These replacements matter for the formal changed regions: the set where two
labellings differ is sensitive to their values on a curve, and if the labelling
reached at the start of an edge were carried through the operations unchanged,
the curve $x=1$ would enter the changed region of a main death and the estimate
would fail there. The estimates are therefore not claimed for every labelling
that agrees with the schedule's off these curves. The replacements do not matter
for the source: the labellings represent the same partition into open chambers,
and no lattice site lies on the curves. A cell center is never on the boundary
of an aligned square of integer side, and never on the curves $x=0$ or $x=1$ of
the band of one of its edges: there $d$ is a half-odd integer, so $d\ne0$ and
$1000d$ is an integer, while $1000\,w(s)=\min(s,n-s)$ is a half-odd integer
(\leanid{TNLean.PEPS.Approximation.IsCellCenter.mem_openSquare_or_notMem_closedSquare},
\leanid{TNLean.PEPS.Approximation.IsCellCenter.notMem_edgeInterface}). Hence the
labels at every lattice site are unchanged by both replacements
(\leanid{TNLean.PEPS.Approximation.blockRepainting_start_site},
\leanid{TNLean.PEPS.Approximation.blockRepainting_final_site}). This is the
source's sampling convention, not a hypothesis of the formal statements.

\end{document}
